\documentclass[a4paper,12pt]{paper}

\usepackage{mycommands}

% Worked-example layout: boxed problem statement followed by a solution.
\newcounter{example}
\renewenvironment{example}[1]{%  % the paper class already defines "example"
  \refstepcounter{example}%
  \begin{mdframed}[linewidth=1pt,innertopmargin=6pt,innerbottommargin=6pt]%
  \noindent\textbf{Example \theexample: #1.}\ }{%
  \end{mdframed}}
\newenvironment{solution}{\par\noindent\textbf{Solution.}\ }{\par\medskip}

\title{Worked examples for Lecture 18: Delay time tomography}
\author{David Al-Attar \\ Michaelmas Term, \the\year}

\begin{document}

\maketitle

\noindent These examples accompany Lecture~18. They are intended as practice in the concrete
calculations that underlie the lecture, and are less demanding than the questions on the
problem sets. Numerical results were checked with the scripts in the \texttt{python}
directory.

\medskip

\noindent Throughout we use a two-block toy model. The reference P-wave velocity is a constant
$\alpha$, and the model consists of two adjacent blocks, each of length $\ell$, in which the
velocity is perturbed by constants $\delta\alpha_{1}$ and $\delta\alpha_{2}$. Rays are straight
lines that pass along the blocks, and so the basis functions of the lecture are the indicator
functions $\varphi_{j}$ of the blocks. The elements of the matrix $\mathbf{A}$ are then
\begin{equation}
    A_{ij} = -\int_{\mathrm{ray}_{i}} \frac{\varphi_{j}}{\alpha^{2}} \dd s = -\frac{\ell_{ij}}{\alpha^{2}},
    \label{eq:1}
\end{equation}
where $\ell_{ij}$ is the length of the $i$th ray within the $j$th block. The factor $-1/\alpha^{2}$ is
common to every element, and it is tidier to take as model parameters the perturbations to the
slowness $1/\alpha$ in each block,
\begin{equation}
    \delta u_{j} = \delta\!\left(\frac{1}{\alpha}\right)_{j} = -\frac{\delta\alpha_{j}}{\alpha^{2}},
    \label{eq:2}
\end{equation}
so that $\mathbf{m}=(\delta u_{1},\delta u_{2})^{T}$ and $A_{ij}=\ell_{ij}$ is simply a path length.
The two descriptions differ only by a rescaling of the columns of $\mathbf{A}$, and the velocity
perturbations can be read off at the end from $\delta\alpha_{j}/\alpha = -\alpha\,\delta u_{j}$.
For numerical values we take $\ell = 100\,\mathrm{km}$ and $\alpha = 8\,\mathrm{km\,s^{-1}}$, so
that the unperturbed travel time along one block is $12.5\,\mathrm{s}$.

\begin{example}{Toy tomography with two blocks}
Consider the following three sets of straight rays through the two-block model.
\begin{enumerate}
    \item[(a)] A single ray that passes along both blocks, with delay time $d_{1}$.
    \item[(b)] The ray of (a) together with a second ray that passes along block 1 only, with delay
    time $d_{2}$.
    \item[(c)] The rays of (b) together with a third ray that passes along block 2 only, with delay
    time $d_{3}$.
\end{enumerate}
In each case write down $\mathbf{A}$ and the equation $\mathbf{d}=\mathbf{A}\mathbf{m}+\mathbf{e}$. Find
the null space of $\mathbf{A}$ in case (a), and solve the problem exactly in case (b). In case (c) the
data $d_{1}=1.2\,\mathrm{s}$, $d_{2}=0.5\,\mathrm{s}$ and $d_{3}=0.4\,\mathrm{s}$ are inconsistent;
find the least squares solution, the predicted data and the residuals, taking the errors to have
equal variances.
\end{example}

\begin{solution}
In case (a) there is one datum and two unknowns, and by eq.~(\ref{eq:1}) with slowness parameters
\begin{equation}
    \mathbf{A} = \begin{pmatrix} \ell & \ell \end{pmatrix},\quad
    d_{1} = \ell\,\delta u_{1} + \ell\,\delta u_{2} + e_{1}.
    \label{eq:3}
\end{equation}
The null space consists of the vectors $\mathbf{m}_{0}$ with $\mathbf{A}\mathbf{m}_{0}=\mathbf{0}$, i.e.\
$\delta u_{1}+\delta u_{2}=0$, and so is the line spanned by $(1,-1)^{T}$. The single ray measures
the sum of the two slowness perturbations and says nothing whatever about their difference: adding
any multiple of $(1,-1)^{T}$ to a solution leaves the datum unchanged, and the problem has no unique
solution even for error-free data. Note that $\mathbf{A}^{T}\mathbf{A} = \ell^{2}\begin{pmatrix} 1 & 1
\\ 1 & 1\end{pmatrix}$ has eigenvectors $(1,1)^{T}$ and $(1,-1)^{T}$ with eigenvalues $2\ell^{2}$
and $0$; the zero eigenvalue is the null space again.

In case (b) the second ray has length $\ell$ in block 1 and zero in block 2, so
\begin{equation}
    \mathbf{A} = \begin{pmatrix} \ell & \ell \\ \ell & 0 \end{pmatrix},\quad
    \det\mathbf{A} = -\ell^{2} \neq 0,\quad
    \mathbf{A}^{-1} = \frac{1}{\ell}\begin{pmatrix} 0 & 1 \\ 1 & -1 \end{pmatrix}.
    \label{eq:4}
\end{equation}
The null space is now trivial and, neglecting errors, the unique solution is
\begin{equation}
    \delta u_{1} = \frac{d_{2}}{\ell},\quad \delta u_{2} = \frac{d_{1}-d_{2}}{\ell}.
    \label{eq:5}
\end{equation}
With $d_{1}=1.2\,\mathrm{s}$ and $d_{2}=0.5\,\mathrm{s}$ this gives $\delta u_{1} =
5\times10^{-3}\,\mathrm{s\,km^{-1}}$ and $\delta u_{2}=7\times10^{-3}\,\mathrm{s\,km^{-1}}$, or
$\delta\alpha_{1}/\alpha = -4.0\,\%$ and $\delta\alpha_{2}/\alpha = -5.6\,\%$: both blocks are slow,
as they must be for positive delays. Because $\mathbf{A}$ is square and invertible the data are fit
exactly whatever their errors, and the errors pass straight into the model,
$\delta\mathbf{m}=\mathbf{A}^{-1}\mathbf{e}$. For uncorrelated errors of variance $\sigma_{d}^{2}$ the
model covariance is
\begin{equation}
    \sigma_{d}^{2}\,\mathbf{A}^{-1}\mathbf{A}^{-T} = \frac{\sigma_{d}^{2}}{\ell^{2}}
    \begin{pmatrix} 1 & -1 \\ -1 & 2 \end{pmatrix},
    \label{eq:6}
\end{equation}
so block 2, which is seen only through the difference $d_{1}-d_{2}$, has twice the variance of block 1,
and the two estimates are anticorrelated.

In case (c) we have three data and two unknowns,
\begin{equation}
    \mathbf{A} = \begin{pmatrix} \ell & \ell \\ \ell & 0 \\ 0 & \ell \end{pmatrix},\quad
    \mathbf{d} = \begin{pmatrix} d_{1} \\ d_{2} \\ d_{3} \end{pmatrix}.
    \label{eq:7}
\end{equation}
Exact solutions exist only if $d_{1}=d_{2}+d_{3}$, and for the given data $d_{2}+d_{3}=0.9\,\mathrm{s}
\neq d_{1}$. With equal variances the covariance matrix is $\mathbf{C}=\sigma_{d}^{2}\mathbf{I}$, the
factor $\sigma_{d}^{-2}$ cancels from the least squares equations of the lecture, and we must solve
$\mathbf{A}^{T}\mathbf{A}\,\mathbf{m} = \mathbf{A}^{T}\mathbf{d}$. Here
\begin{equation}
    \mathbf{A}^{T}\mathbf{A} = \ell^{2}\begin{pmatrix} 2 & 1 \\ 1 & 2 \end{pmatrix},\quad
    (\mathbf{A}^{T}\mathbf{A})^{-1} = \frac{1}{3\ell^{2}}\begin{pmatrix} 2 & -1 \\ -1 & 2 \end{pmatrix},\quad
    \mathbf{A}^{T}\mathbf{d} = \ell\begin{pmatrix} d_{1}+d_{2} \\ d_{1}+d_{3} \end{pmatrix},
    \label{eq:8}
\end{equation}
and hence
\begin{equation}
    \mathbf{m} = \frac{1}{3\ell}\begin{pmatrix} d_{1}+2d_{2}-d_{3} \\ d_{1}-d_{2}+2d_{3} \end{pmatrix}
    = \begin{pmatrix} 6\times10^{-3} \\ 5\times10^{-3} \end{pmatrix}\mathrm{s\,km^{-1}},
    \label{eq:9}
\end{equation}
which corresponds to $\delta\alpha_{1}/\alpha=-4.8\,\%$ and $\delta\alpha_{2}/\alpha=-4.0\,\%$. The
predicted data are $\mathbf{A}\mathbf{m} = (1.1, 0.6, 0.5)^{T}\,\mathrm{s}$, and the residuals are
\begin{equation}
    \mathbf{A}\mathbf{m}-\mathbf{d} = -\frac{d_{1}-d_{2}-d_{3}}{3}\begin{pmatrix} 1 \\ -1 \\ -1 \end{pmatrix}
    = \begin{pmatrix} -0.1 \\ 0.1 \\ 0.1 \end{pmatrix}\mathrm{s}.
    \label{eq:10}
\end{equation}
The residual vector is orthogonal to both columns of $\mathbf{A}$, as it must be at a least squares
minimum, and is proportional to $(1,-1,-1)^{T}$: the only combination of the data that the model
cannot reproduce is the inconsistency $d_{1}-d_{2}-d_{3}$, and least squares distributes it equally
between the three measurements. The example shows the three regimes met in practice: an
underdetermined problem with a null space, an exactly determined one in which noise is amplified,
and an overdetermined one whose data can be fit only in the least squares sense.
\end{solution}

\begin{example}{Regularised least squares in the toy problem}
Return to case (a) of Example~1, with a single datum $d$, and measure delay times in units of their
standard deviation so that $\mathbf{C}=\mathbf{I}$. Taking $\mathbf{B}=\mathbf{I}$, compute the
regularised least squares solution
\begin{equation}
    \mathbf{m}(\lambda) = (\mathbf{A}^{T}\mathbf{A}+\lambda\mathbf{I})^{-1}\mathbf{A}^{T}d
    \label{eq:11}
\end{equation}
explicitly, and discuss its limits as $\lambda\to0$ and $\lambda\to\infty$. Find the data misfit
$\|\mathbf{A}\mathbf{m}-d\|$ and the model norm $\|\mathbf{m}\|$ as functions of $\lambda$, and
hence the trade-off curve between them. Evaluate the results for $d=1.2\,\mathrm{s}$ and
$\lambda/\ell^{2}=0$, $0.5$, $2$ and $8$.
\end{example}

\begin{solution}
With $\mathbf{A}=(\ell\;\;\ell)$ we have
\begin{equation}
    \mathbf{A}^{T}\mathbf{A}+\lambda\mathbf{I} = \begin{pmatrix} \ell^{2}+\lambda & \ell^{2} \\
    \ell^{2} & \ell^{2}+\lambda \end{pmatrix},\quad
    \det(\mathbf{A}^{T}\mathbf{A}+\lambda\mathbf{I}) = (\ell^{2}+\lambda)^{2}-\ell^{4} = \lambda(2\ell^{2}+\lambda),
    \label{eq:12}
\end{equation}
which is positive for every $\lambda>0$, in agreement with the general argument in the lecture that
the regularised normal matrix is positive definite. Inverting the two-by-two matrix,
\begin{equation}
    (\mathbf{A}^{T}\mathbf{A}+\lambda\mathbf{I})^{-1} = \frac{1}{\lambda(2\ell^{2}+\lambda)}
    \begin{pmatrix} \ell^{2}+\lambda & -\ell^{2} \\ -\ell^{2} & \ell^{2}+\lambda \end{pmatrix},
    \label{eq:13}
\end{equation}
and since $\mathbf{A}^{T}d = \ell d\,(1,1)^{T}$ we find
\begin{equation}
    \mathbf{m}(\lambda) = \frac{\ell d}{\lambda(2\ell^{2}+\lambda)}
    \begin{pmatrix} \lambda \\ \lambda \end{pmatrix}
    = \frac{\ell d}{2\ell^{2}+\lambda}\begin{pmatrix} 1 \\ 1 \end{pmatrix}.
    \label{eq:14}
\end{equation}
The same result follows more quickly from the eigenvectors noted in Example~1. The matrix
$\mathbf{A}^{T}\mathbf{A}+\lambda\mathbf{I}$ has eigenvalue $2\ell^{2}+\lambda$ on $(1,1)^{T}$ and
$\lambda$ on $(1,-1)^{T}$; the right-hand side $\mathbf{A}^{T}d$ lies entirely along $(1,1)^{T}$, and
so is simply divided by $2\ell^{2}+\lambda$. In particular the regularised solution never acquires a
component along the null-space direction $(1,-1)^{T}$: with $\mathbf{B}=\mathbf{I}$ and
$\mathbf{m}_{0}=\mathbf{0}$, damping selects the solution that puts nothing into the null space.

As $\lambda\to0$, $\mathbf{m}\to(d/2\ell)(1,1)^{T}$. This fits the datum exactly, since
$\ell(d/2\ell + d/2\ell) = d$, and among all exact solutions $(d/2\ell)(1,1)^{T}+a(1,-1)^{T}$ it has the
smallest norm, $\|\mathbf{m}\|^{2} = d^{2}/2\ell^{2} + 2a^{2}$ being minimised at $a=0$. Thus the small
damping limit is the minimum-norm solution, which attributes the delay equally to the two blocks. As
$\lambda\to\infty$, $\mathbf{m}\to\mathbf{0}$: the regularisation term dominates and the model is
pulled entirely to the reference model, whatever the data say.

For the misfit and norm we have $\mathbf{A}\mathbf{m} = 2\ell^{2}d/(2\ell^{2}+\lambda)$, and so
\begin{equation}
    \|\mathbf{A}\mathbf{m}-d\| = \frac{\lambda\,d}{2\ell^{2}+\lambda},\quad
    \|\mathbf{m}\| = \frac{\sqrt{2}\,\ell\,d}{2\ell^{2}+\lambda}.
    \label{eq:15}
\end{equation}
The misfit rises monotonically from $0$ to $d$ as $\lambda$ increases, while the norm falls
monotonically from $d/\sqrt{2}\ell$ to $0$. Eliminating $\lambda$ between the two expressions,
\begin{equation}
    \|\mathbf{A}\mathbf{m}-d\| = d - \sqrt{2}\,\ell\,\|\mathbf{m}\|,
    \label{eq:16}
\end{equation}
so the trade-off curve in the (norm, misfit) plane is a straight line joining the minimum-norm
solution $(d/\sqrt{2}\ell,0)$ to the reference model $(0,d)$. For $d=1.2\,\mathrm{s}$ and
$\ell=100\,\mathrm{km}$ the numbers are as follows.
\begin{center}
\begin{tabular}{cccc}
    $\lambda/\ell^{2}$ & $\delta u_{1}=\delta u_{2}$ ($\mathrm{s\,km^{-1}}$) &
    $\|\mathbf{A}\mathbf{m}-d\|$ (s) & $\|\mathbf{m}\|$ ($\mathrm{s\,km^{-1}}$) \\[2pt]
    \hline\\[-8pt]
    $0$   & $6.0\times10^{-3}$ & $0$    & $8.49\times10^{-3}$ \\
    $0.5$ & $4.8\times10^{-3}$ & $0.24$ & $6.79\times10^{-3}$ \\
    $2$   & $3.0\times10^{-3}$ & $0.60$ & $4.24\times10^{-3}$ \\
    $8$   & $1.2\times10^{-3}$ & $0.96$ & $1.70\times10^{-3}$ \\
\end{tabular}
\end{center}
The value $\lambda=2\ell^{2}$, equal to the single non-zero eigenvalue of $\mathbf{A}^{T}\mathbf{A}$,
halves the amplitude of the solution and misfits the datum by half its value. The trade-off curve
here is a straight line because $\mathbf{A}$ has rank one; the characteristic bend, or ``corner'', of
trade-off curves in real problems arises only when the singular values of $\mathbf{A}$ span a range,
different components of the model then being switched off at different values of $\lambda$.
\end{solution}

\begin{example}{Fermat's principle checked exactly}
A medium of uniform P-wave velocity $\alpha$ contains a slab $0<x<h$ in which the velocity is
$\alpha(1+\epsilon)$. A source at $(-a,-b)$ and a receiver at $(h+a,b)$ lie on opposite sides of
the slab, so that the unperturbed ray is a straight line crossing the slab obliquely at the angle
$\theta$ to its normal, where $\tan\theta = 2b/(2a+h)$. Show from Snell's law that the exact travel
time is given by
\begin{align}
    T(\epsilon) &= \frac{2a}{\alpha\cos\theta_{1}} + \frac{h}{\alpha(1+\epsilon)\cos\theta_{2}}, \nonumber \\
    \sin\theta_{2} &= (1+\epsilon)\sin\theta_{1},\quad 2a\tan\theta_{1} + h\tan\theta_{2} = 2b.
    \label{eq:17}
\end{align}
Show that the linearised delay time of the lecture is $\delta T = -\epsilon\ell_{\mathrm{in}}/\alpha$,
with $\ell_{\mathrm{in}}=h/\cos\theta$ the length of the unperturbed ray inside the slab, that the
first-order term of $T(\epsilon)$ agrees with this, and that the contribution of ray bending to the
delay time is $O(\epsilon^{2})$ and negative. Confirm this numerically for $\alpha=8\,\mathrm{km\,s^{-1}}$,
$h=100\,\mathrm{km}$, $a=200\,\mathrm{km}$, $b=250\,\mathrm{km}$ and
$\epsilon = 0.01$, $0.03$ and $0.1$.
\end{example}

\begin{solution}
Within each homogeneous region the perturbed ray is straight, and at the two plane interfaces it
refracts according to Snell's law, the horizontal slowness $p=\sin\theta_{1}/\alpha
= \sin\theta_{2}/\alpha(1+\epsilon)$ being conserved. Let $\theta_{1}$ be the angle to the slab normal
of the two outer segments and $\theta_{2}$ that of the segment inside the slab. The outer segments
each have length $a/\cos\theta_{1}$ and the inner segment $h/\cos\theta_{2}$, which gives the travel
time in eq.~(\ref{eq:17}); the transverse distances covered, $a\tan\theta_{1}$ twice and
$h\tan\theta_{2}$ once, must sum to $2b$, which is the final condition. For given $\epsilon$ these two
conditions determine $\theta_{1}$ through a single transcendental equation, which is readily solved
numerically. At $\epsilon=0$ we have $\theta_{1}=\theta_{2}=\theta$ and the ray is the straight line
with $T_{0} = (2a+h)/\alpha\cos\theta$.

The velocity perturbation is $\delta v = \epsilon\alpha$ inside the slab and zero outside, and the
linearised delay time of the lecture is the integral along the \emph{unperturbed} ray,
\begin{equation}
    \delta T = -\int_{\mathrm{ray}}\frac{\delta v}{v^{2}}\dd s = -\frac{\epsilon\alpha}{\alpha^{2}}\,\ell_{\mathrm{in}}
    = -\frac{\epsilon\,\ell_{\mathrm{in}}}{\alpha},\quad \ell_{\mathrm{in}} = \frac{h}{\cos\theta}.
    \label{eq:18}
\end{equation}

To compare with the exact result it is clearest to make the Fermat argument explicit. Any
three-segment path is fixed by its entry point $(0,y_{1})$ and exit point $(h,y_{2})$, and has travel
time
\begin{align}
    T(\epsilon;\mathbf{y}) &= T_{\mathrm{out}}(\mathbf{y}) + \frac{T_{\mathrm{in}}(\mathbf{y})}{1+\epsilon}, \nonumber \\
    T_{\mathrm{out}} &= \frac{\sqrt{a^{2}+(y_{1}+b)^{2}}+\sqrt{a^{2}+(b-y_{2})^{2}}}{\alpha},\quad
    T_{\mathrm{in}} = \frac{\sqrt{h^{2}+(y_{2}-y_{1})^{2}}}{\alpha},
    \label{eq:19}
\end{align}
with $\mathbf{y}=(y_{1},y_{2})^{T}$. Snell's law is precisely the condition
$\nabla_{\mathbf{y}}T(\epsilon;\mathbf{y})=\mathbf{0}$, and the true ray, $\mathbf{y}(\epsilon)$, is the
minimiser of $T(\epsilon;\cdot)$; in particular the unperturbed ray $\mathbf{y}_{0}$ satisfies
$\nabla T_{\mathrm{out}}(\mathbf{y}_{0}) = -\nabla T_{\mathrm{in}}(\mathbf{y}_{0}) \equiv -\mathbf{g}$.
Writing $\mathbf{y}(\epsilon) = \mathbf{y}_{0}+\delta\mathbf{y}$ and expanding,
\begin{equation}
    T(\epsilon) = T(\epsilon;\mathbf{y}_{0}) + \nabla_{\mathbf{y}}T(\epsilon;\mathbf{y}_{0})\cdot\delta\mathbf{y}
    + \frac{1}{2}\delta\mathbf{y}^{T}\mathbf{H}\,\delta\mathbf{y} + \cdots,
    \label{eq:20}
\end{equation}
where $\mathbf{H}$ is the Hessian of $T_{0}(\mathbf{y}) = T(0;\mathbf{y})$ at $\mathbf{y}_{0}$. The first
term is the travel time through the perturbed slab along the unperturbed path,
\begin{equation}
    T(\epsilon;\mathbf{y}_{0}) = T_{0} + \left(\frac{1}{1+\epsilon}-1\right)\frac{\ell_{\mathrm{in}}}{\alpha}
    = T_{0} - \frac{\epsilon\,\ell_{\mathrm{in}}}{\alpha} + \frac{\epsilon^{2}\ell_{\mathrm{in}}}{\alpha} - \cdots,
    \label{eq:21}
\end{equation}
whose first-order part is exactly eq.~(\ref{eq:18}). In the second term of eq.~(\ref{eq:20}) the
gradient is $\nabla_{\mathbf{y}}T(\epsilon;\mathbf{y}_{0}) = -\epsilon\mathbf{g}/(1+\epsilon)$, which
vanishes at $\epsilon=0$ by Fermat's principle, while $\delta\mathbf{y}=O(\epsilon)$ because
the ray moves continuously with the model. The ray-bending contribution to the travel time is
therefore $O(\epsilon^{2})$: the first-order delay time is obtained by integrating along the old
ray, as claimed in the lecture. We can go one step further. Stationarity of the travel time at
$\mathbf{y}_{0}+\delta\mathbf{y}$ gives $\mathbf{H}\,\delta\mathbf{y} = \epsilon\mathbf{g}+O(\epsilon^{2})$, and
substituting into eq.~(\ref{eq:20}),
\begin{equation}
    T(\epsilon) - T(\epsilon;\mathbf{y}_{0}) = -\epsilon^{2}\mathbf{g}^{T}\mathbf{H}^{-1}\mathbf{g}
    + \frac{1}{2}\epsilon^{2}\mathbf{g}^{T}\mathbf{H}^{-1}\mathbf{g} + O(\epsilon^{3})
    = -\frac{1}{2}\epsilon^{2}\,\mathbf{g}^{T}\mathbf{H}^{-1}\mathbf{g} + O(\epsilon^{3}),
    \label{eq:22}
\end{equation}
which is negative because $\mathbf{H}$ is positive definite at a minimum. Bending always shortens
the travel time relative to the frozen path, whatever the sign of $\epsilon$, since the true ray is
the fastest path in the perturbed medium. For the present geometry the quantities are elementary. A
straight segment of length $L$ making angle $\theta$ with the $x$-axis has travel time
$\sqrt{\Delta x^{2}+\Delta y^{2}}/\alpha$ whose second derivative with respect to an end-point
ordinate is $\cos^{2}\theta/\alpha L$, and so at $\mathbf{y}_{0}$
\begin{equation}
    \mathbf{H} = \frac{\cos^{3}\theta}{\alpha}\begin{pmatrix} 1/a+1/h & -1/h \\ -1/h & 1/a+1/h \end{pmatrix},\quad
    \mathbf{g} = \nabla T_{\mathrm{in}} = \frac{\sin\theta}{\alpha}\begin{pmatrix} -1 \\ 1 \end{pmatrix}.
    \label{eq:23}
\end{equation}
The vector $\mathbf{g}$ is an eigenvector of $\mathbf{H}$ with eigenvalue $(\cos^{3}\theta/\alpha)(1/a+2/h)$,
and hence
\begin{equation}
    -\frac{1}{2}\mathbf{g}^{T}\mathbf{H}^{-1}\mathbf{g} = -\frac{\sin^{2}\theta\,ah}{\alpha\cos^{3}\theta\,(h+2a)}
    = -\frac{\ell_{\mathrm{in}}}{\alpha}\,\frac{a\tan^{2}\theta}{h+2a}.
    \label{eq:24}
\end{equation}
Collecting eqs.~(\ref{eq:21}), (\ref{eq:22}) and (\ref{eq:24}), the exact delay time is
\begin{equation}
    T(\epsilon)-T_{0} = -\frac{\epsilon\,\ell_{\mathrm{in}}}{\alpha}
    + \frac{\epsilon^{2}\ell_{\mathrm{in}}}{\alpha}\left(1 - \frac{a\tan^{2}\theta}{h+2a}\right) + O(\epsilon^{3}).
    \label{eq:25}
\end{equation}
The two second-order pieces have distinct origins: the first is the non-linearity of the slowness
$1/v$ in $\epsilon$ and would be absent had we perturbed the slowness rather than the velocity; the
second is ray bending, and vanishes at normal incidence where the ray is not bent at all.

For the numerical values given, $\tan\theta = 500/500$, so $\theta = 45^{\circ}$,
$\ell_{\mathrm{in}}=h\sqrt{2}=141.4\,\mathrm{km}$, $T_{0}=500\sqrt{2}/8 = 88.39\,\mathrm{s}$ and
$\ell_{\mathrm{in}}/\alpha = 17.68\,\mathrm{s}$. The bending coefficient in eq.~(\ref{eq:24}) is
$-17.68\times0.4 = -7.07\,\mathrm{s}$, and the total second-order coefficient in eq.~(\ref{eq:25}) is
$17.68\times0.6=10.61\,\mathrm{s}$. Solving eq.~(\ref{eq:17}) numerically we obtain the following.
\begin{center}
\begin{tabular}{ccccc}
    $\epsilon$ & $-\epsilon\ell_{\mathrm{in}}/\alpha$ (s) & $T(\epsilon)-T_{0}$ (s) &
    error (ms) & bending part (ms) \\[2pt]
    \hline\\[-8pt]
    $0.01$ & $-0.17678$ & $-0.17573$ & $1.05$  & $-0.70$ \\
    $0.03$ & $-0.53033$ & $-0.52114$ & $9.19$  & $-6.25$ \\
    $0.1$  & $-1.76777$ & $-1.67393$ & $93.8$  & $-66.9$ \\
\end{tabular}
\end{center}
Here ``error'' is the difference between the exact and the linearised delay times, and ``bending
part'' is $T(\epsilon)-T(\epsilon;\mathbf{y}_{0})$, the part of that difference due to the movement of
the ray. Increasing $\epsilon$ by a factor of three multiplies the error by $8.8$ and by a factor of
ten multiplies it by $90$, close to the factors of $9$ and $100$ expected for a quadratic, and the
second-order formula $10.61\,\epsilon^{2}\,\mathrm{s}$ predicts $1.06$, $9.5$ and $106\,\mathrm{ms}$,
the discrepancies being the $O(\epsilon^{3})$ terms. Even for a $10\,\%$ velocity anomaly, far larger
than anything in the mantle, the linearised delay time is in error by only $6\,\%$, and the
ray-bending part of that error is smaller in magnitude than, and opposite in sign to, the part due to
the non-linearity of the slowness in $\epsilon$. This is why the linear relation
$\mathbf{d}=\mathbf{A}\mathbf{m}$ is adequate for global tomography.
\end{solution}

\begin{example}{Resolution and the Bayesian view}
Return once more to case (a) of Example~1, $\mathbf{A}=(\ell\;\;\ell)$, and suppose the datum has
variance $\sigma_{d}^{2}$ and that the prior on the model is Gaussian with mean $\mathbf{m}_{0}=\mathbf{0}$
and covariance $\mathbf{C}_{m}=\sigma_{m}^{2}\mathbf{I}$. Compute the posterior mean $\mathbf{m}_{p}$ and
covariance $\mathbf{C}_{p}$ explicitly. Show that the posterior variance is reduced only along the
direction $(1,1)^{T}$ and is unchanged along the null-space direction $(1,-1)^{T}$, and relate
the result to the regularised solution of Example~2. Evaluate everything for
$d=1.2\,\mathrm{s}$, $\sigma_{d}=0.1\,\mathrm{s}$ and $\sigma_{m}=2\times10^{-3}\,\mathrm{s\,km^{-1}}$.
\end{example}

\begin{solution}
Let $\mathbf{P}_{\parallel}$ and $\mathbf{P}_{\perp}$ be the orthogonal projectors onto the unit
vectors $(1,1)^{T}/\sqrt{2}$ and $(1,-1)^{T}/\sqrt{2}$,
\begin{equation}
    \mathbf{P}_{\parallel} = \frac{1}{2}\begin{pmatrix} 1 & 1 \\ 1 & 1 \end{pmatrix},\quad
    \mathbf{P}_{\perp} = \frac{1}{2}\begin{pmatrix} 1 & -1 \\ -1 & 1 \end{pmatrix},\quad
    \mathbf{P}_{\parallel}+\mathbf{P}_{\perp} = \mathbf{I}.
    \label{eq:26}
\end{equation}
From Example~1, $\mathbf{A}^{T}\mathbf{A} = 2\ell^{2}\mathbf{P}_{\parallel}$, and so the inverse posterior
covariance of the lecture is
\begin{equation}
    \mathbf{C}_{p}^{-1} = \frac{\mathbf{A}^{T}\mathbf{A}}{\sigma_{d}^{2}} + \frac{\mathbf{I}}{\sigma_{m}^{2}}
    = \left(\frac{2\ell^{2}}{\sigma_{d}^{2}}+\frac{1}{\sigma_{m}^{2}}\right)\mathbf{P}_{\parallel}
    + \frac{1}{\sigma_{m}^{2}}\mathbf{P}_{\perp}.
    \label{eq:27}
\end{equation}
A matrix of the form $\mu_{\parallel}\mathbf{P}_{\parallel}+\mu_{\perp}\mathbf{P}_{\perp}$ is inverted by
inverting its two eigenvalues, and hence
\begin{equation}
    \mathbf{C}_{p} = c\,\mathbf{P}_{\parallel} + \sigma_{m}^{2}\mathbf{P}_{\perp},\quad
    c = \frac{\sigma_{m}^{2}\sigma_{d}^{2}}{\sigma_{d}^{2}+2\ell^{2}\sigma_{m}^{2}} < \sigma_{m}^{2}.
    \label{eq:28}
\end{equation}
The posterior variance of the component of $\mathbf{m}$ along $(1,1)^{T}/\sqrt{2}$ is $c$, which is
smaller than the prior variance $\sigma_{m}^{2}$ by the factor $\sigma_{d}^{2}/(\sigma_{d}^{2}+2\ell^{2}\sigma_{m}^{2})$,
while that along $(1,-1)^{T}/\sqrt{2}$ is $\sigma_{m}^{2}$, exactly as before the datum was collected.
The data reduce our uncertainty only in the directions they are sensitive to; in the null space the
posterior is the prior, and any apparent knowledge of the difference between the blocks in the
final model comes entirely from the prior. Written out,
\begin{equation}
    \mathbf{C}_{p} = \frac{1}{2}\begin{pmatrix} \sigma_{m}^{2}+c & c-\sigma_{m}^{2} \\
    c-\sigma_{m}^{2} & \sigma_{m}^{2}+c \end{pmatrix},
    \label{eq:29}
\end{equation}
so the two blocks are negatively correlated a posteriori, with correlation coefficient
$(c-\sigma_{m}^{2})/(c+\sigma_{m}^{2})$: the datum fixes the sum of the perturbations, and so if one
block is faster than the posterior mean the other is probably slower. In the limit of a perfect
datum, $\sigma_{d}\to0$, we have $c\to0$, $\mathbf{C}_{p}\to\sigma_{m}^{2}\mathbf{P}_{\perp}$, and the
correlation tends to $-1$.

The posterior mean is $\mathbf{m}_{p} = \mathbf{C}_{p}\mathbf{A}^{T}d/\sigma_{d}^{2}$ since $\mathbf{m}_{0}=\mathbf{0}$,
and as $\mathbf{A}^{T}d = \ell d\,(1,1)^{T}$ lies along $(1,1)^{T}$,
\begin{equation}
    \mathbf{m}_{p} = \frac{c\,\ell d}{\sigma_{d}^{2}}\begin{pmatrix} 1 \\ 1 \end{pmatrix}
    = \frac{\ell d}{2\ell^{2}+\sigma_{d}^{2}/\sigma_{m}^{2}}\begin{pmatrix} 1 \\ 1 \end{pmatrix}.
    \label{eq:30}
\end{equation}
Comparing with eq.~(\ref{eq:14}), this is the regularised least squares solution of Example~2 with
\begin{equation}
    \lambda = \frac{\sigma_{d}^{2}}{\sigma_{m}^{2}},
    \label{eq:31}
\end{equation}
which is the lecture's identification $\lambda\mathbf{B}=\mathbf{C}_{m}^{-1}$ once the data variance,
set to unity in Example~2, is restored: with $\mathbf{C}=\sigma_{d}^{2}\mathbf{I}$ the regularised
normal matrix is $\mathbf{A}^{T}\mathbf{A}/\sigma_{d}^{2}+\lambda\mathbf{B}$, and
$\lambda\mathbf{B}=\mathbf{I}/\sigma_{m}^{2}$ gives $\mathbf{A}^{T}\mathbf{A}+(\sigma_{d}^{2}/\sigma_{m}^{2})\mathbf{I}$.
The damping parameter is thus the ratio of how noisy we think the data are to how large we think the
model is: noisier data or a tighter prior both mean heavier damping.

It is also instructive to ask how the posterior mean is related to the true model. If the datum
was generated as $d = \mathbf{A}\mathbf{m}_{\mathrm{true}}+e$, then
\begin{equation}
    \mathbf{m}_{p} = \mathbf{R}\,\mathbf{m}_{\mathrm{true}} + \frac{\mathbf{C}_{p}\mathbf{A}^{T}e}{\sigma_{d}^{2}},\quad
    \mathbf{R} = \frac{\mathbf{C}_{p}\mathbf{A}^{T}\mathbf{A}}{\sigma_{d}^{2}}
    = \frac{2\ell^{2}c}{\sigma_{d}^{2}}\mathbf{P}_{\parallel}
    = \frac{2\ell^{2}\sigma_{m}^{2}}{\sigma_{d}^{2}+2\ell^{2}\sigma_{m}^{2}}\mathbf{P}_{\parallel},
    \label{eq:32}
\end{equation}
where $\mathbf{R}$ is called the resolution matrix. The recovered model is a filtered version of the
truth: the mean of the two blocks is recovered with amplitude slightly less than one, and the
difference between them is not recovered at all.

For the given numbers, $2\ell^{2}\sigma_{m}^{2} = 0.08\,\mathrm{s^{2}}$ and $\sigma_{d}^{2}=0.01\,\mathrm{s^{2}}$,
so $c=\sigma_{m}^{2}/9$: the posterior standard deviation along $(1,1)^{T}/\sqrt{2}$ is
$\sigma_{m}/3 = 0.67\times10^{-3}\,\mathrm{s\,km^{-1}}$, and along $(1,-1)^{T}/\sqrt{2}$ it remains
$2\times10^{-3}\,\mathrm{s\,km^{-1}}$. Explicitly,
\begin{equation}
    \mathbf{C}_{p} = \begin{pmatrix} 2.22 & -1.78 \\ -1.78 & 2.22 \end{pmatrix}\times10^{-6}\,\mathrm{s^{2}\,km^{-2}},
    \label{eq:33}
\end{equation}
with marginal standard deviations $1.49\times10^{-3}\,\mathrm{s\,km^{-1}}$ for each block, a
correlation coefficient of $-0.8$, and $\mathbf{R} = (8/9)\mathbf{P}_{\parallel}$. The damping parameter
is $\lambda = \sigma_{d}^{2}/\sigma_{m}^{2} = 2500\,\mathrm{km^{2}}$, i.e.\ $\lambda/\ell^{2}=0.25$, and
the posterior mean is $\mathbf{m}_{p} = (5.33, 5.33)\times10^{-3}\,\mathrm{s\,km^{-1}}$, corresponding to
$\delta\alpha/\alpha=-4.3\,\%$ in each block. Note that each block's marginal uncertainty is reduced
only modestly, from $2.0$ to $1.5\times10^{-3}\,\mathrm{s\,km^{-1}}$, even though their sum is known
three times better than a priori: marginal error bars on individual parameters can badly
misrepresent what has been learned when the posterior is strongly correlated.
\end{solution}

\begin{example}{Sensitivity along the ray}
A ray passes through a homogeneous medium of velocity $v$. The velocity is perturbed by a small
constant $\delta$ inside a sphere of radius $R$ whose centre lies on the ray. Evaluate the
linearised delay time, and comment on how it depends on $R$. What happens when the centre of the
sphere is displaced a distance $b$ perpendicular to the ray? Compare with the finite-frequency
sensitivity kernels of Lecture~19.
\end{example}

\begin{solution}
The linearised delay time is the integral along the unperturbed ray,
\begin{equation}
    \delta T = -\int_{\mathrm{ray}}\frac{\delta v}{v^{2}}\dd s.
    \label{eq:34}
\end{equation}
The ray is a straight line through the centre of the sphere, so it intersects the sphere along a
diameter, of length $2R$, on which $\delta v=\delta$, and vanishes elsewhere. Hence
\begin{equation}
    \delta T = -\frac{2R\,\delta}{v^{2}},
    \label{eq:35}
\end{equation}
a fast anomaly ($\delta>0$) advancing the arrival and a slow one delaying it, as it should. For
$v=8\,\mathrm{km\,s^{-1}}$, $R=50\,\mathrm{km}$ and a $1\,\%$ slow anomaly, $\delta=-0.08\,\mathrm{km\,s^{-1}}$,
we obtain $\delta T = 0.125\,\mathrm{s}$, a typical size for a teleseismic P-wave delay.

The delay time is proportional to $R$, the length of the ray inside the anomaly, and not to the
volume $\tfrac{4}{3}\pi R^{3}$ of the anomaly: halving $R$ halves $\delta T$ but reduces the volume
eightfold. This is the concrete meaning of the statement in the lecture that the functional
derivative $DT(v)$ is a generalised function concentrated on the ray. If $DT(v)$ were an ordinary
bounded function $K(\mathbf{x})$, then $\langle DT(v)|\delta v\rangle = \delta\int_{\mathrm{sphere}}K\dd^{3}\mathbf{x}$
would be $O(R^{3})$ as $R\to0$, not $O(R)$. In the language of the block parameterisation, the
sphere is a single basis function and $-2R/v^{2}$ is the corresponding element of $\mathbf{A}$.

If the centre is displaced a distance $b$ from the ray, the ray cuts the sphere along a chord of
length $2\sqrt{R^{2}-b^{2}}$ for $b<R$, and misses it entirely for $b>R$, so
\begin{equation}
    \delta T = \begin{cases} -\dfrac{2\sqrt{R^{2}-b^{2}}\,\delta}{v^{2}}, & b<R, \\[8pt] 0, & b\geq R. \end{cases}
    \label{eq:36}
\end{equation}
For the numbers above and $b=30\,\mathrm{km}$ the delay is $0.100\,\mathrm{s}$, and for
$b=80\,\mathrm{km}$ it is exactly zero. Ray theory therefore predicts that an anomaly is invisible to
a ray that does not pass through it, however close it comes, and that the delay is largest when the
ray passes through the centre. The finite-frequency kernels of Lecture~19 differ on both counts. They
are non-zero throughout the first Fresnel zone, so an anomaly just off the ray still produces a
delay, and they vanish on the ray itself, because a scatterer on the ray produces no phase delay
relative to the direct wave. The two descriptions agree only when the anomaly is large compared with
the width of the Fresnel zone, in which case the integral of the finite-frequency kernel over the
anomaly reduces to the chord length and eq.~(\ref{eq:36}) is recovered. For a $50\,\mathrm{km}$
anomaly and a $1\,\mathrm{Hz}$ teleseismic P-wave, whose Fresnel zone in the mantle is a few hundred
kilometres wide, ray theory is already a rough guide.
\end{solution}

\end{document}
